\subsection{The retained period constant is an effective deformation parameter}
\label{sec:retained-period-deformation}

Keep the functions, lattice marking, base coordinate and local fillings of
\cref{sec:candidate-geometry}. In particular, fix the actual cusp-regular
solution $\beta^{\rm part}$, not a replacement with a cusp pole, and write
\[
 \beta_c=\beta^{\rm part}+c,\qquad
 \mathcal P=\{c\in\mathbb C:\operatorname{Im}c<-M\},
 \quad M=\max_{B\setminus\{p_0\}}D_{\rm part}.
 \tag{PD1}
\]
The maximum and the exact admissible domain were proved in
\cref{prop:period-exact-admissible-constants}. The following calculation
uses the analytic construction, not its integral cohomology or sphere
recognition. Its purpose is to determine whether the freedom in $c$
actually changes the constructed complex manifold.

\begin{theorem}[Exact period deformation]
\label{thm:retained-period-deformation}
For every open disc $Q$ with $\overline Q\subset\mathcal P$, the retained
construction gives a proper holomorphic submersion
\[
 \pi:\mathcal X_Q\longrightarrow Q,
 \qquad F:\mathcal X_Q\longrightarrow B\times Q,
 \qquad \pi=\operatorname{pr}_Q\circ F,
 \tag{PD2}
\]
whose fibre over $c$ is the analytically constructed threefold $X_c$ and
whose restriction there is the original $f_c:X_c\to B$.
The three distinguished values remain the original
$t(p_1)=0$, $t(p_2)=1$, $t(p_0)=\infty$.

Fix a regular base point and one of its original lifts $z$, and fix
$c_*\in Q$. With all functions below evaluated at that $z$, set
\[
 T=\operatorname{Im}\tau>0,\quad m=\operatorname{Im}\mu,\quad
 D=\operatorname{Im}\beta_{c_*}-\frac{6m^2}{T}<0,\quad a=\frac mT,
 \qquad P=\begin{pmatrix}0&0\\-a&1\end{pmatrix}.
 \tag{PD3}
\]
For $\delta=c-c_*$, the exact marked real-linear map from the universal
cover of the fibre at $c_*$ to that at $c$ is
\[
 w=A_\delta\zeta+B_\delta\overline\zeta,\qquad
 A_\delta=I+\frac{\delta}{2iD}P,\qquad
 B_\delta=-\frac{\delta}{2iD}P.
 \tag{PD4}
\]
It maps the original period lattice onto the new period lattice, and
therefore descends to a real-analytic marked diffeomorphism of the two
tori. In the explicit convention that the transported $(0,1)$ tangent
space is spanned by
$\partial_{\overline\zeta_k}+\sum_j\nu^j_k\partial_{\zeta_j}$,
its exact coefficient matrix and its first derivative are
\[
 \nu_\delta=\frac{\delta}{2iD+\delta}P,\qquad
 \dot\nu_z=\frac{1}{2iD}\,
 \partial_{\zeta_2}\otimes
 (d\overline\zeta_2-a\,d\overline\zeta_1).
 \tag{PD5}
\]
The class of $\dot\nu_z$ in
$H^{0,1}(f_{c_*}^{-1}(t(z)),T^{1,0})$ is nonzero.
With the tangent-exact-sequence convention
$\operatorname{KS}(\partial_c)=[\bar\partial h]$ for a smooth
$(1,0)$ lift $h$ of $\partial_c$, the actual fibre class is
$\operatorname{KS}_z(\partial_c)=-[\dot\nu_z]$.
Moreover the first-order deformation of the entire $X_{c_*}$ obtained
from (PD2) is nontrivial: its Kodaira--Spencer class in
$H^1(X_{c_*},T_{X_{c_*}})$ is nonzero. In particular
$h^1(X_{c_*},T_{X_{c_*}})\geq1$.
\end{theorem}

\begin{proof}
\emph{Construction with the parameter retained.}
On an ordinary lifted base open set use the original period map
\[
 \Pi(z,c)=
 \begin{pmatrix}
 6\mu(z)&\tau(z)&1&0\\
 \beta^{\rm part}(z)+c&\mu(z)&0&1
 \end{pmatrix}:
 \Lambda\otimes\mathbb R\longrightarrow\mathbb C^2.
 \tag{PD6}
\]
Its real determinant is $\operatorname{Im}\tau\,D_c$, which never
vanishes on $Q$. Thus $(z,c,u)\mapsto[z,c,\Pi(z,c)u]$ identifies the
quotient, as a real-analytic family, with the original marked real torus
over this base open set times $Q$. The lattice action is holomorphic in
$(z,c,\zeta)$; it is free and properly discontinuous locally uniformly
on compact base-parameter sets, since the inverse real period matrices
are uniformly bounded there. The quotient is therefore a complex
manifold. Projection to $Q$ is a submersion because it is so before
the covering quotient. Its map to the lifted base times $Q$ is proper:
over any compact subset the image of that subset times a fixed closed
lattice parallelotope is the full inverse image.

The deck matrices are the original integral $M_g$ and the original
matrices $R_g(z)$, with
$\Pi(gz,c)=R_g(z)\Pi(z,c)M_g$. The $R_g$ do not depend on $c$.
Consequently the ordinary base quotient is jointly holomorphic.
Proper discontinuity there follows from proper discontinuity on the
base, exactly as in \cref{lem:middle-period-proper}, with compact
parameter sets carried along.

At a finite filling retain $m_1=3$, $m_2=4$,
$s_j(g_jz)=e^{-2\pi i/m_j}s_j(z)$, and the actual logarithmic lift
\[
 (z,c,\zeta)\longmapsto
 \left(g_jz,c,R_{g_j}(z)\zeta+
                  \Pi(g_jz,c)\frac{v_j}{m_j}\right).
 \tag{PD7}
\]
In the real marking this is $(z,c,u)\mapsto
(g_jz,c,A_ju+v_j/m_j)$, independent of $c$.
Every nonidentity power remains free: the invariant integral covector
$\gamma$ takes the translation of its $k$-th power to $k/3$ or $-k/4$,
which is not an integer for $0<k<m_j$.
The finite quotient is therefore a local biholomorphism and is a
submersion to $Q$. The same compact-parallelotope argument proves
properness over $D_j\times Q$. Its map is still exactly
$t_j=s_j^{m_j}$. On the punctured overlap the translating section is
\[
 \frac{\log s_j(z)}{2\pi i}\Pi(z,c)v_j.
 \tag{PD8}
\]
A change of logarithm adds an actual period. It is therefore a
branch-independent holomorphic torus section, jointly in $z,c$.
The calculation (2.9e) applies with $c$ retained and proves the
punctured-overlap conjugacy, not merely equality of the monodromy orders.

For the cusp retain the same infinite unimodular fan and the full matrix
\[
 C(t_c,c)=
 \begin{pmatrix}6\mu(t_c)&h(t_c)\\
 b^{\rm part}(t_c)-h(t_c)+c&\mu(t_c)
 \end{pmatrix},\qquad
 B_0=\begin{pmatrix}0&1\\-1&0\end{pmatrix}.
 \tag{PD9}
\]
Choose a closed cusp disc of radius $r_1$ on which these functions are
holomorphic. Put
\[
 K_Q=\max_{|t_c|\le r_1,\ c\in\overline Q}
          \|-2\pi\operatorname{Im}C(t_c,c)\|,
 \qquad 0<\epsilon<\min(r_1,1),\quad
          |\log\epsilon|>2K_Q.
 \tag{PD10}
\]
The action on $Y_\epsilon\times Q$ is the original one,
\[
 (x,t_c,c)\longmapsto
 \left(e^{2\pi iC(t_c,c)\lambda}
                  t_c^{B_0\lambda}x,t_c,c\right),
 \quad\lambda\in\mathbb Z^2.
 \tag{PD11}
\]
It is holomorphic across the boundary by the same toric automorphism
and nowhere-zero torus multiplier. The estimates in (2.8) now hold
uniformly on $Q$:
\[
 B_{t_c,c}=B_0+
       \frac{-2\pi\operatorname{Im}C(t_c,c)}{\log|t_c|},\quad
 \|B_{t_c,c}-B_0\|<\tfrac12,\quad
 \|B_{t_c,c}\lambda\|>\tfrac12\|\lambda\|\quad(\lambda\ne0).
 \tag{PD12}
\]
They exclude stabilizers on the dense torus and bound the finite set
of translations that can join each pair of bounded charts, uniformly
in $c$. At $t_c=0$ a nonzero translation moves the bounded cell, so
there is no stabilizer there either. A compact set is contained in finitely
many of the open charts $\mathcal U_\sigma(2,\epsilon)\times Q$ from
\cref{lem:toric-chart-control}.
The chart-pair estimate proves proper discontinuity on the whole
product, including sequences tending to its boundary.

For clarity properness is also uniform, not only pointwise in $c$.
For $\epsilon_1<\epsilon$, let $K_0$ be the finite union of closed
unit polydiscs associated to triangles meeting the closed radius
$\sqrt2$ ball, intersected with $|t_c|\le\epsilon_1$, exactly as in
\cref{thm:candidate-cusp-geometry}. For every $c$ and $t_c\ne0$,
choose $\lambda$ so that
$\|B_{t_c,c}^{-1}y+\lambda\|_\infty\le1/2$.
The translated logarithmic position has norm at most
$3/(2\sqrt2)<\sqrt2$, and hence lies in a chart defining $K_0$.
At $t_c=0$, the six closed bidiscs covering $E_0$ give the same
representatives after the exact translation by $B_0v$, with multiplier
$\exp(2\pi iC(0,c)B_0v)$ retained. Therefore, for compact $K\subset Q$,
the image of the compact $K_0\times K$ is the full inverse image of
$\{|t_c|\le\epsilon_1\}\times K$. Every compact set in the cusp
base times $Q$ is closed in a set of this form. This proves properness.
The quotient is a submersion to $Q$ because its covering charts are
open subsets of $Y_\epsilon\times Q$.

Exponentiating the original logarithmic torus coordinates gives the
jointly holomorphic cusp overlap, with period identity
$Z=sB_0+C(t_c,c)$. Glue it and (PD8) to the middle family.
The special base discs are disjoint. Each resulting inverse image of
a member of this base cover is its entire Hausdorff local piece.
Two points with different $(t,c)$ values are separated on the base;
two with the same value lie in one Hausdorff open piece and can be
separated there. This proves Hausdorffness. Local complex charts give
a complex fourfold and the submersion $\pi$.
Properness of $F$ follows by covering any compact set of $B\times Q$
by finitely many relatively compact sets inside these base pieces and
using the just proved local properness. Since $B$ is compact, projection
$B\times Q\to Q$ is proper, proving (PD2). The fibres and overlaps
are precisely those in the original analytic construction.

\emph{The exact variation on a regular fibre.}
Write $u=(u_1,u_2,u_3,u_4)^t$ in the original ordered lattice basis and
$\zeta=\Pi(z,c_*)u$. Taking imaginary parts gives the two exact equations
\[
 \operatorname{Im}\zeta_1=6m u_1+T u_2,\qquad
 \operatorname{Im}\zeta_2=(D+6m^2/T)u_1+m u_2.
 \tag{PD13}
\]
Subtracting $a$ times the first equation from the second yields
\[
 u_1=\frac{\operatorname{Im}\zeta_2-a\operatorname{Im}\zeta_1}{D}.
 \tag{PD14}
\]
In the changed periods the same $u$ maps to
$w=\Pi(z,c)u=\zeta+\delta e_2u_1$. Substitution of
$\operatorname{Im}\zeta=(\zeta-\overline\zeta)/(2i)$ proves (PD4).
This also proves its asserted lattice map and real invertibility:
it is exactly $\Pi(z,c)\Pi(z,c_*)^{-1}$ as a real-linear map.
Its real determinant is $(D+\operatorname{Im}\delta)/D>0$.
Neither the nonholomorphic part nor the original marking is dropped.

Here $P^2=P$, and
\[
 \det A_\delta=1+\frac\delta{2iD},\qquad
 A_\delta^{-1}=I-\frac{\delta}{2iD+\delta}P.
 \tag{PD15}
\]
The denominator can vanish only for $\delta=-2iD$. At that parameter
$D+\operatorname{Im}\delta=-D>0$, which is outside the admissible
domain. Thus this inverse is defined at every $c\in Q$.
Annihilating $w$ by the declared new $(0,1)$ fields gives
$A_\delta\nu_\delta+B_\delta=0$, which proves (PD5), including
its sign. All its coefficients are constant in the fibre coordinates.
Thus its first derivative is a globally defined $\bar\partial$-closed
tangent-valued $(0,1)$ form on the torus.

It cannot be $\bar\partial V$ for a globally defined smooth vector
field $V$. Indeed the invariant holomorphic frame
$\partial_{\zeta_1},\partial_{\zeta_2}$ writes $V$ with periodic
smooth coefficients. The integral of each of their derivatives
$\partial_{\overline\zeta_k}V^j$ against the translation-invariant
probability volume is zero: real constant vector fields preserve this
volume, and integration of their derivatives on the compact torus is
zero. But the coefficient of
$\partial_{\zeta_2}\otimes d\overline\zeta_2$ in $\dot\nu_z$ is
the nonzero constant $1/(2iD)$. Its integral is that same constant.
This contradiction proves nonzero Dolbeault class without a dimension
count or an appeal to numerical period variation.

The conversion to the tangent-exact-sequence convention retains a
minus sign. In the moving coordinates the real-marking lift has
$(1,0)$ part $h=\partial_c+u_1\partial_{w_2}$ at $c_*$.
Since (PD14) gives
\[
 \bar\partial u_1=-\frac{d\overline\zeta_2-a\,d\overline\zeta_1}{2iD},
\]
the connecting class $[\bar\partial h]$ is exactly $-[\dot\nu_z]$.
Both conventions therefore give a nonzero class; they are not identified
without this sign conversion.

\emph{The intrinsic pencil in the parameter family.}
Let $\mathcal S_2\subset\mathcal X_Q$ be the reduced finite divisor
over $p_2\times Q$, so $F^*(p_2\times Q)=4\mathcal S_2$.
The proof of \cref{thm:canonical-ring-recovery} works for relative
forms over $Q$, as can be checked on each original chart.
On the middle family the frame is $d\zeta_1\wedge d\zeta_2$ relative
to $B\times Q$. At the cusp the relative canonical frame over $Q$ is
$(dx_1/x_1)\wedge(dx_2/x_2)\wedge dt_c$.
Extra $dc$ terms in transition differentials vanish in relative forms.
The finite determinant characters $\chi_j$, their correcting units
$e^{\varphi_j(s_j)}$, and the modular function $\Theta$ used in
(9.4)--(9.6) are independent of $c$, since $R_g,\tau,\mu$ are.
Thus their coefficient identities hold jointly over $Q$.
The evaluation exponents remain $m_j-1-a_j=0,2$.
It follows by the actual same transition functions and divisors that
\[
 \omega_{\mathcal X_Q/Q}\simeq\mathcal O_{\mathcal X_Q}(-2\mathcal S_2),
 \qquad
 \omega_{\mathcal X_Q/Q}^{-2}\simeq
          F^*\operatorname{pr}_B^*\mathcal O_B(p_2).
 \tag{PD16}
\]
Here the base identifications retain $h=1/(t-1)$ and
$\eta_B=dt/(t-1)^2$ from (CR1); these functions also do not depend on $c$.
In particular (PD16) is a relative line-bundle isomorphism, not a
collection of separately chosen isomorphisms of its fibres.

We need its complete degree-two system after an infinitesimal base
change, and prove this explicitly. Put
$A=\mathbb C[\epsilon]/(\epsilon^2)$, base-change by
$c=c_*+\epsilon$, write $X_A$ for the resulting deformation and
$L_A=\omega_{X_A/A}^{-2}$. For $L_*=\omega_{X_{c_*}}^{-2}$,
flatness of the deformation and local freeness of $L_A$ give
\[
 0\longrightarrow L_*\xrightarrow{\ \epsilon\ }L_A
        \longrightarrow L_*\longrightarrow0.
 \tag{PD17}
\]
The two sections of $\mathcal O_B(p_2)$ used in (CR4), namely
$\sigma_{p_2}$ and $\ell=A_0/(t-1)+B_0'$ with $A_0\ne0$,
pull back by (PD16) to sections $s_0,s_1$ of $L_A$ whose reductions
are a basis of $H^0(X_{c_*},L_*)$. This basis assertion is exactly
the proved degree-two case of (CR2)--(CR4).
Any section of $L_A$ reduces to a complex linear combination of these
two reductions; subtract that lifted combination. Exactness of (PD17)
then writes the remainder as $\epsilon$ times a section of $L_*$,
again a combination of that basis. Conversely, if an $A$-linear
combination of $s_0,s_1$ is zero, reduce modulo $\epsilon$ and then
use the injective first arrow of (PD17) to see that all four complex
coefficients vanish. Therefore
\[
 H^0(X_A,L_A)=A s_0\oplus A s_1.
 \tag{PD18}
\]
The system is basepoint-free and its morphism is $f_A$ followed by
the fixed coordinate map $t\mapsto[t-1:A_0+B_0'(t-1)]$.
This proves the required base-change statement without presuming that
pointwise dimensions alone imply it.

\emph{Nontriviality for the complete threefold.}
Suppose its first-order deformation were trivial, with an isomorphism
$X_A\simeq X_{c_*}\times\operatorname{Spec}A$ reducing to the identity.
Intrinsic relative canonical bundles and their complete spaces of
sections are preserved by that isomorphism. By (PD18), it would
therefore identify $f_A$ with $a_A\circ f_{c_*}$ for a projective
linear change $a_A\in\operatorname{PGL}_2(A)$ reducing to the identity.

This change fixes all three distinguished points even to first order.
Here is a scheme-sensitive check. In the finite covering charts
$f_A$ is $t_j=s_j^{m_j}$, so its relative critical ideal is
$(s_j^{m_j-1})$, and its restriction to that critical scheme has
$t_j=0$. The map from the parameter ring $A$ to this critical
scheme is injective; it contains a constant-parameter section given
by $s_j=0$ and the zero section of the covering torus.
At the cusp, in a chart containing a triple point,
$t_c=z_0z_1z_2$ and the critical ideal is
$(z_1z_2,z_0z_2,z_0z_1)$. Again $t_c=0$ in its ring and $A$ injects,
using $z_0=z_1=z_2=0$. At a double chart the same assertion follows
from $t_c=z_0z_1$; a smooth chart has no critical point.
These critical ideals are intrinsic to the differential of the map
over $A$, and the charts are local biholomorphisms onto $X_A$.
Away from the three distinguished base values the map is a
submersion, so it has no critical point.
Consequently the scheme-theoretic critical-value sections are exactly
the three constant sections $p_1,p_2,p_0$. An automorphism reducing
to the identity cannot permute these three distinct closed points.
Its value at each is therefore the same point over $A$, not a
nonzero infinitesimal displacement: such a displacement would be a
nonzero element of $\epsilon\mathbb C$ in an injective copy of $A$
on the just displayed critical chart.

For completeness write $a_A(t)=t+\epsilon(b_0+b_1t+b_2t^2)$.
Fixing $0$ gives $b_0=0$, fixing $1$ gives $b_0+b_1+b_2=0$,
and fixing infinity, in $u=1/t$, gives $b_2=0$.
Hence $a_A$ is the identity. The hypothesized trivialization thus
preserves $f$, and restricting it to any regular base value would
trivialize the torus deformation (PD4). Its nonzero class (PD5)
excludes this.

Finally, the zero Kodaira--Spencer class is precisely the condition for
such a first-order trivialization. One may see the needed direction
directly: choose local product charts for the holomorphic submersion
$\pi$ and differentiate their transition functions at $c_*$. The
resulting tangent-valued overlap functions form a Cech $1$-cocycle.
If its class vanishes, after refinement it is the coboundary of local
holomorphic vector fields. Changing each chart by that vector field
times $\epsilon$ cancels the first-order overlap terms, producing
the identity-reducing trivialization. This is the usual
Kodaira--Spencer connecting class (for the standard deformation-space
interpretation see \cite[Section~67, p.~15]{McMullen1996Complex});
the explicitly recorded conversion
from the plus-graph convention in (PD5) changes its sign, not its
vanishing. The contradiction proves the final assertion.
\end{proof}

The theorem establishes a genuinely varying analytic family with its
original marked fibration and its intrinsic degree-two canonical pencil.
It does not assert that every two distinct parameters give nonisomorphic
unmarked manifolds, or that this family exhausts their deformation space.
The global integral topology and the CDP argument are not inputs to this
deformation calculation.


The global identifications within this retained family are computed in
\cref{pshift:thm:classification}: they are precisely $c'-c\in2\mathbb Z$,
including the finite-filling obstruction and every cusp overlap.
